
\documentclass[a4paper,12pt]{article}

% Zeichenkodierung und Sprache
\usepackage[utf8]{inputenc}
\usepackage[T1]{fontenc}
\usepackage[ngerman]{babel}
\usepackage{lmodern}

% Seitenlayout
\usepackage{geometry}
\geometry{margin=2.5cm}

% Abbildungen, Tabellen und Zeichnungen
\usepackage{graphicx}
\usepackage{float}
\usepackage{booktabs}
\usepackage{tikz}
\usetikzlibrary{arrows.meta,positioning}

% Quellcode und Pfadangaben
\usepackage{listings}
\usepackage{xcolor}
\usepackage{url}

% Einstellungen für Befehlsblöcke
\lstset{
    basicstyle=\ttfamily\small,
    frame=single,
    breaklines=true,
    columns=fullflexible,
    keepspaces=true,
    showstringspaces=false,
    captionpos=b
}

% Absatzformatierung
\setlength{\parindent}{0pt}
\setlength{\parskip}{0.5em}

\title{Bericht}
\author{Jiaqi Wang}
\date{\today}

\begin{document}

\maketitle

\section{ROS 2-Steuerung einer WS2812-LED über WLAN}

In diesem Versuch wurde die integrierte WS2812-LED des ESP32-S3 über ROS~2 und WLAN gesteuert. Der ESP32-S3 arbeitet als WLAN-Access-Point und startet gleichzeitig einen UDP-Server. Nach dem Verbinden des Computers mit dem WLAN des ESP32-S3 können UDP-Nachrichten direkt an den Mikrocontroller gesendet werden.

Auf der ROS~2-Seite wurde ein Python-Knoten implementiert. Dieser abonniert das Topic \texttt{/led\_color} vom Typ \texttt{std\_msgs/msg/String}. Die empfangenen RGB-Werte werden überprüft und anschließend als UDP-Paket an die IP-Adresse \texttt{192.168.4.1} auf Port \texttt{8080} übertragen.

Der ESP32-S3 empfängt Nachrichten im Format \texttt{R,G,B}, beispielsweise \texttt{255,0,0}. Nach erfolgreichem Empfang werden die drei Farbwerte ausgelesen und an die Bibliothek \texttt{Adafruit NeoPixel} übergeben, wodurch die integrierte WS2812-LED unmittelbar ihre Farbe ändert.

Durch die Kombination aus ROS~2, Python und WLAN konnte eine zuverlässige Kommunikation zwischen PC und ESP32-S3 realisiert werden. Das Verfahren eignet sich als Grundlage für zukünftige ROS~2-basierte Anwendungen, bei denen Sensoren oder Aktoren drahtlos über WLAN gesteuert werden.


\section{Entwurf der Grundplatte mit Inkscape}

Um den Schrittmotor, die Steuerplatine und den Akku zu einer kompakten Einheit zusammenzufassen, wurde eine Grundplatte mit der Vektorgrafiksoftware \textit{Inkscape} entworfen. Ziel war es, alle Komponenten sicher auf einer gemeinsamen Trägerplatte zu befestigen und den späteren Aufbau zu vereinfachen.

Zunächst wurden mit einer Schieblehre die Positionen der Befestigungsbohrungen vermessen. Insgesamt wurden vier Bohrungen vorgesehen. Für den Schrittmotor wurden zwei Bohrungen mit einem Durchmesser von \textbf{M5} erstellt, während für die Steuerplatine zwei Bohrungen mit einem Durchmesser von \textbf{M3} vorgesehen wurden.

Die endgültige Größe der Grundplatte beträgt \textbf{31\,cm $\times$ 20\,cm}. Um die Anforderungen der Laserschneidmaschine zu erfüllen, wurden sämtliche Konturen ausschließlich aus roten Linien mit einer Linienbreite von \textbf{0,03\,mm} erstellt. Alle Rechtecke und Kreise besitzen keine Flächenfüllung (\textit{Fill: None}), sodass ausschließlich die Schneidkonturen exportiert werden.

Nach Abschluss des Entwurfs wurde die Zeichnung im SVG-Format unter dem Dateinamen \texttt{base\_v10.svg} gespeichert und anschließend für den Laserschneidprozess verwendet.

\end{document}

